\section{Discussion: how models get good at specific tasks}\label{sec:discussion}

The 2026 pattern --- one frontier, several winners --- is the predictable
output of how frontier models are now built. We see at least four levers.

\paragraph{1. Domain-weighted post-training.} Reinforcement learning on
verifiable, domain-specific rewards (unit tests for code, proof checkers for
math, task-completion signals for agents) lets a lab buy capability in a chosen
vertical. Moonshot's agentic suite results and OpenAI's terminal/agentic lead
both track heavy RL investment in those task
families~\cite{tomshardware2026k3,openai2026gpt56}.

\paragraph{2. Architecture as targeting.} Kimi~K3's hybrid linear attention
(Kimi Delta Attention) and expert-routing changes were chosen for long-context
agentic efficiency, and the payoff appeared precisely on long-horizon
benchmarks (BrowseComp 91.2, state of the art at
release)~\cite{venturebeat2026k3,tomshardware2026k3}.

\paragraph{3. Inference-time structure.} GPT-5.6's Ultra mode makes
multi-agent fan-out a first-class feature, buying Terminal-Bench points
(88.8\%~$\to$~91.9\%) at ${\sim}3\times$ cost~\cite{vellum2026tiers,eesel2026gpt56}.
Our self-consistency experiment (Section~\ref{sec:selfcons}) is the same lever
at toy scale: capability is partly a function of inference-time orchestration.

\paragraph{4. Benchmark-shaped optimization (and its limits).} When a benchmark
becomes a target, labs optimize toward it, and two failure modes follow: the
benchmark saturates (GLUE, MMLU, soon SWE-bench Verified), or gains fail to
transfer off-distribution. The Witness evaluation of Opus~5 is the cleanest
2026 example: a ${\sim}4\times$ record on ARC-AGI-3, but only a statistical tie
with K3 and Fable~5 on held-out puzzles of the same genre, including one
environment where Opus~5 stated the hidden rules before its first move ---
suggesting familiarity, not just reasoning~\cite{decoder2026opus5,techtimes2026opus5}.
METR's maintainer-review study adds a second caution from a different angle:
a growing share of SWE-bench-passing patches would be rejected by human
maintainers (quality, breakage, or missing the issue's point), and the
grader--maintainer gap is widening over model
generations~\cite{metr2026swebench}. Targeting works, and that is exactly why
single-number claims deserve suspicion.

\paragraph{Implication for evaluation.} Score-versus-cost curves, held-out
transfer suites, and maintainer-grade review are becoming the minimum standard
of evidence. Vendors already report curves~\cite{valletta2026opus5}; ARC Prize
now administers runs independently with published
traces~\cite{arcprize2026opus5, techtimes2026opus5}; and METR's maintainer
study is a template for auditing what benchmark passes actually mean.

\paragraph{Implication for buyers.} The economic shape of 2026 --- a budget
tier at 4\% of flagship cost that matches last quarter's flagship on most
professional work, plus a fragmented frontier --- means the unit of procurement
is no longer ``a model'' but a routing policy with per-task effort settings.
The 2023 question ``which model is best?'' has been replaced by ``which
combination, at which effort, for which task?''

\paragraph{Limitations.} This study has several. (i) Mid-2026 scores mix
vendor and independent evaluations; we label provenance but cannot fully
deconfound harness differences (e.g., GPT-5.5 scores 85.1\% vendor-reported
vs.\ 82.6\% reconstructed from Vals AI's independent
harness)~\cite{localaimaster2026gpt55,vals2026swebench}. (ii) The SWE-bench
trend fit pools vendor and independent numbers and spans only 22 months; the
$5.8\times$ figure should be read as an order of magnitude. (iii) The routing
oracle uses benchmark identities as a proxy for task identity; real routers
must infer task type from unlabeled traffic. (iv) The self-consistency
experiment uses a 0.5B model on grade-school math by design --- it
demonstrates the mechanism cheaply and reproducibly, not the frontier effect
size. (v) The field moves weekly; figures are current as of August 2026.
